\section{Conclusion}

This retrospective study evaluated BESS arbitrage recommendations as decisions rather than forecasts. A common LP/oracle contour converted candidate information into feasible schedules and measured regret against realized-price optimal schedules. The main slice contains 18 dates crossed with five configured BESS profiles sharing one DAM price process. Raw and horizon-calibrated NBEATSx are co-primary sensitivity sources: V2+ mean regret is 193.36 UAH for the raw source and 174.77 UAH for the calibrated source, versus 310.58 UAH for strict and 225.44/206.37 UAH for V2. V2+ has lower mean regret than both comparators in all four adjacent rolling windows for both sources. These are descriptive retrospective estimates, not confirmatory holdout results.

A post-defense lineage audit materially narrows the learned-challenger claims. The artifact historically named \code{dt_v2_plus} is a random forest, not a Decision Transformer. Its 360 training candidates exactly mirror its 360 evaluation candidates after a timestamp shift, and all four nonfallback profile rows occur on one delivery date. Its 168.16 UAH value is therefore an in-packet construction diagnostic, not evidence of temporal generalization. Separate Hugging Face transformer values use the mirrored packet, while a 32-day read-model audit has no realized-regret outcome. None is promoted.

The contribution is a claim-safe architecture and evidence-governance pattern, not a new forecasting, optimization, random-forest, or sequence-model method. Official rows are used for published targets; forecasts supply scenarios for unpublished horizons; schedules share one evaluator; challengers can abstain; and the operator surface remains a read model. The repository exposes code and a frozen evidence subset \cite{fefelov2026repo}.
